% ═══════════════════════════════════════════════════════════════
% MT-01a-ML: Musterlösung Minitest Begrüßungen
% Spanisch Klasse 7 - Sequenz 1: Empezamos
% ═══════════════════════════════════════════════════════════════

\documentclass[11pt, a4paper]{article}

\usepackage[utf8]{inputenc}
\usepackage[T1]{fontenc}
\usepackage[ngerman]{babel}
\usepackage{helvet}
\renewcommand{\familydefault}{\sfdefault}

\usepackage[margin=2cm]{geometry}
\usepackage{enumitem}
\usepackage{tabularx}
\usepackage{booktabs}
\usepackage{xcolor}
\usepackage{fancyhdr}
\usepackage{lastpage}
\usepackage{amssymb}

% FARBEN
\definecolor{spanisch}{HTML}{c41e3a}
\definecolor{grau}{HTML}{888888}
\definecolor{hellgrau}{HTML}{f5f5f5}
\definecolor{loesung}{HTML}{c62828}

\pagestyle{fancy}
\fancyhf{}
\renewcommand{\headrulewidth}{0.4pt}
\renewcommand{\footrulewidth}{0.4pt}

\lhead{Spanisch -- Klasse 7}
\chead{\textbf{MT-01a -- MUSTERLÖSUNG}}
\rhead{\today}

\lfoot{10 Minuten}
\cfoot{}
\rfoot{Seite \thepage\ von \pageref{LastPage}}

\newcommand{\punktebox}[1]{\hfill\fbox{\small \_\_\_/#1}}
\newcommand{\luecke}[1]{\rule{#1}{0.4pt}}
\newcommand{\es}[1]{\textcolor{spanisch}{\textbf{#1}}}
\newcommand{\lsg}[1]{\textcolor{loesung}{\textbf{#1}}}

\begin{document}

% ═══════════════════════════════════════════════════════════════
% KOPFBEREICH
% ═══════════════════════════════════════════════════════════════

\begin{center}
    {\Large\bfseries\textcolor{spanisch}{Minitest: Begrüßungen}}

    \vspace{5pt}
    {\large\textcolor{loesung}{-- MUSTERLÖSUNG --}}
\end{center}

\vspace{10pt}

% ═══════════════════════════════════════════════════════════════
% AUFGABE 1: Zuordnung (6 BE)
% ═══════════════════════════════════════════════════════════════

\noindent\textbf{\textcolor{spanisch}{Aufgabe 1 -- Begrüßungen zuordnen}} \punktebox{6}

\vspace{10pt}

\begin{enumerate}[label=\arabic*.]
    \item \es{¡Hola!} = \lsg{b) Hallo!}
    \item \es{¡Buenos días!} = \lsg{a) Guten Morgen!}
    \item \es{¡Buenas noches!} = \lsg{c) Guten Abend!}
    \item \es{¡Adiós!} = \lsg{d) Auf Wiedersehen!}
    \item \es{¡Hasta luego!} = \lsg{e) Bis später!}
    \item \es{¡Hasta mañana!} = \lsg{f) Bis morgen!}
\end{enumerate}

\textit{Je 1 BE pro richtige Zuordnung.}

\vspace{15pt}

% ═══════════════════════════════════════════════════════════════
% AUFGABE 2: Tageszeit (6 BE)
% ═══════════════════════════════════════════════════════════════

\noindent\textbf{\textcolor{spanisch}{Aufgabe 2 -- Welche Begrüßung passt?}} \punktebox{6}

\vspace{10pt}

\begin{tabularx}{\textwidth}{|l|X|}
\hline
\textbf{Uhrzeit} & \textbf{Begrüßung} \\
\hline
8:00 Uhr morgens & \lsg{Buenos días} \\
\hline
15:00 Uhr nachmittags & \lsg{Buenas tardes} \\
\hline
21:00 Uhr abends & \lsg{Buenas noches} \\
\hline
10:00 Uhr vormittags & \lsg{Buenos días} \\
\hline
17:00 Uhr nachmittags & \lsg{Buenas tardes} \\
\hline
23:00 Uhr nachts & \lsg{Buenas noches} \\
\hline
\end{tabularx}

\vspace{5pt}
\textit{Je 1 BE pro richtige Begrüßung. Zeiten: morgens/vormittags (bis 14 Uhr) = días, nachmittags (14-20 Uhr) = tardes, abends/nachts (ab 20 Uhr) = noches.}

\vspace{15pt}

% ═══════════════════════════════════════════════════════════════
% AUFGABE 3: Grammatik (6 BE)
% ═══════════════════════════════════════════════════════════════

\noindent\textbf{\textcolor{spanisch}{Aufgabe 3 -- Buenos oder Buenas?}} \punktebox{6}

\vspace{10pt}

\begin{enumerate}[label=\alph*)]
    \item \lsg{Buenos} días \hfill \textit{(1 BE)}
    \item \lsg{Buenas} tardes \hfill \textit{(1 BE)}
    \item \lsg{Buenas} noches \hfill \textit{(1 BE)}
\end{enumerate}

\vspace{10pt}

\noindent \textbf{Begründung (3 BE):}

\vspace{5pt}

\noindent
\lsg{$\boxtimes$} Weil ``día'' maskulin und ``tarde/noche'' feminin ist \\
$\square$ Weil ``día'' kürzer ist \\
$\square$ Weil es morgens ist

\vspace{5pt}
\textit{Erklärung: Das Adjektiv ``bueno/buena'' passt sich dem Geschlecht des Substantivs an. ``El día'' ist maskulin (Ausnahme: endet auf -a, ist aber maskulin!), ``la tarde'' und ``la noche'' sind feminin.}

\vspace{15pt}

% ═══════════════════════════════════════════════════════════════
% AUFGABE 4: Verabschiedung (4 BE)
% ═══════════════════════════════════════════════════════════════

\noindent\textbf{\textcolor{spanisch}{Aufgabe 4 -- Verabschiedungen ergänzen}} \punktebox{4}

\vspace{10pt}

\begin{enumerate}[label=\alph*)]
    \item Bis bald! = ¡Hasta \lsg{pronto}! \hfill \textit{(2 BE)}
    \item Bis später! = ¡Hasta \lsg{luego}! \hfill \textit{(2 BE)}
\end{enumerate}

\vspace{15pt}

% ═══════════════════════════════════════════════════════════════
% AUFGABE 5: Satzzeichen (4 BE)
% ═══════════════════════════════════════════════════════════════

\noindent\textbf{\textcolor{spanisch}{Aufgabe 5 -- Spanische Satzzeichen}} \punktebox{4}

\vspace{10pt}

\noindent
$\square$ Sie stehen nur am Ende des Satzes \\
\lsg{$\boxtimes$} Sie stehen am Anfang UND am Ende des Satzes \hfill \textit{(2 BE)} \\
$\square$ Sie werden nicht verwendet

\vspace{10pt}

\noindent ``Hallo, wie geht's?'' auf Spanisch:

\vspace{5pt}

\noindent \lsg{¡Hola! ¿Qué tal?} \hfill \textit{(2 BE)}

\vspace{5pt}
\textit{Auch akzeptiert: ``¡Hola! ¿Cómo estás?'' oder ähnliche korrekte Varianten. Wichtig: Umgekehrte Zeichen am Satzanfang!}

% ═══════════════════════════════════════════════════════════════
% PUNKTEÜBERSICHT
% ═══════════════════════════════════════════════════════════════

\vspace{20pt}
\hrule
\vspace{10pt}

\noindent
\begin{tabularx}{\textwidth}{|X|c|c|c|c|c||c|}
\hline
\textbf{Aufgabe} & 1 & 2 & 3 & 4 & 5 & \textbf{Gesamt} \\
\hline
\textbf{max. Punkte} & 6 & 6 & 6 & 4 & 4 & \textbf{26} \\
\hline
\end{tabularx}

\vspace{15pt}

\noindent\textbf{Notenskala (14-NP):}

\vspace{5pt}

\begin{tabular}{|c|c|c|c|c|c|c|c|c|c|c|c|c|c|c|}
\hline
NP & 14 & 13 & 12 & 11 & 10 & 9 & 8 & 7 & 6 & 5 & 4 & 3 & 2 & 1 \\
\hline
\% & 100 & 96 & 91 & 86 & 80 & 73 & 67 & 60 & 53 & 47 & 40 & 33 & 27 & 20 \\
\hline
BE & 26 & 25 & 24 & 22 & 21 & 19 & 17 & 16 & 14 & 12 & 10 & 9 & 7 & 5 \\
\hline
\end{tabular}

\end{document}
